\documentclass{article}

% Language setting
\usepackage[utf8]{inputenc}
\usepackage[english]{babel}

% Page size and margins
\usepackage[a4paper,top=2cm,bottom=2cm,left=3cm,right=3cm]{geometry}

% Useful packages
\usepackage{amsmath}
\usepackage{amssymb}
\usepackage{graphicx}
\usepackage[colorlinks=true, allcolors=blue]{hyperref}
\usepackage{listings}
\usepackage{xcolor}
\usepackage{float}
\usepackage{booktabs}

\lstdefinestyle{pythonstyle}{
    language=Python,
    backgroundcolor=\color{gray!10},
    keywordstyle=\color{blue}\bfseries,
    commentstyle=\color{green!50!black},
    stringstyle=\color{orange},
    numberstyle=\tiny\color{gray},
    basicstyle=\ttfamily\small,
    breaklines=true,
    showstringspaces=false,
    frame=single,
    numbers=left,
    tabsize=2
}

\title{Real-Time Payment Fraud Detection for an E-Commerce Platform\\[0.3em]
\large Business Analytics Group Project --- From Raw Data to a Deployable Model}

\author{%
  \textbf{Team:} [Insert team member names] \\[0.4em]
  \textbf{Course:} Business Analytics \\[0.4em]
  \textbf{Source:} \url{https://github.com/Trungdn4-458311/FraudDetectionBE}%
}
\date{}

\begin{document}
\maketitle

% =========================
\begin{abstract}
Payment fraud causes substantial financial loss and erodes customer
trust on e-commerce platforms. Acting as the risk-analytics function of
a platform, this project answers a concrete business question:
\emph{which transactions are likely fraudulent, and how should the
platform balance blocking fraud against approving legitimate orders
without unnecessary friction?} We combine \textbf{two data sources}: the
simulated mobile-money dataset \textbf{PaySim} ($\sim$6.36M
transactions, fraud rate $0.129\%$) and \textbf{9 contextual features}
generated in Python with Faker (account age, new-device flag,
shipping/billing mismatch, IP distance, etc.). The pipeline covers
exploratory data analysis (EDA), documented before/after cleaning,
feature engineering, and training of \textbf{three classifiers}
(Logistic Regression, Random Forest, XGBoost) over \textbf{two feature
spaces} (Base/Full), evaluated with metrics suited to imbalanced data
(PR-AUC, Recall, F1) plus a \textbf{cost-based metric} for selecting an
operating threshold. A key methodological point: the 9 synthetic
features are generated \emph{conditioned on the label}, introducing
\textbf{label leakage}; we quantify its effect through the Base-vs-Full
comparison. Our best models (XGBoost/Random Forest) reach PR-AUC
$>0.99$, but we show that the near-perfect scores largely reflect the
\emph{easiness} of the data (PaySim is nearly separable via balance
features) rather than genuine real-world generalization --- an
essential conclusion for an honest report.
\end{abstract}

\noindent\textbf{Keywords:} fraud detection; class imbalance; PaySim;
synthetic data; label leakage; PR-AUC; cost-based threshold; XGBoost;
Random Forest.

% =========================
\section{Introduction}
\label{sec:intro}

An e-commerce platform processes thousands of transactions per minute,
and even a fraud rate well under $5\%$ translates into large absolute
losses. We frame the task as \textbf{binary classification} on a
\textbf{severely imbalanced} dataset: flag high-risk transactions in
real time while limiting the number of legitimate customers blocked.

\subsection{Business KPIs}
\label{subsec:kpi}
We track three indicators:
\begin{itemize}
    \item \textbf{Fraud rate} --- the proportion of fraudulent
    transactions (measures imbalance and risk scale).
    \item \textbf{False-positive rate (FPR)} --- the proportion of
    legitimate transactions wrongly blocked (measures customer friction).
    \item \textbf{Financial loss avoided} --- the fraud amount stopped
    by the model (measures business value).
\end{itemize}
Across the full dataset the fraud rate is $0.129\%$ ($8{,}213$ out of
$6{,}362{,}620$ transactions), with a \textbf{total money at risk} of
$\sim$$1.21\times 10^{10}$ units and a mean of $\sim$$1.47$M per
fraudulent transaction. This is the baseline against which
Section~\ref{subsec:cost} is compared.

\subsection{Contributions and structure}
The report follows the first five modules of the project:
Section~\ref{sec:data} (Module~1) describes the data and synthetic
generation; Section~\ref{sec:eda} (Module~2) explores the data;
Section~\ref{sec:clean} (Module~3) cleans it;
Section~\ref{sec:feat} (Module~4) engineers features;
Section~\ref{sec:model} (Module~5) trains, evaluates, and selects an
operating threshold.

% =========================
\section{Data and Synthetic Generation (Module 1)}
\label{sec:data}

\subsection{Source 1 --- PaySim}
PaySim is a simulated mobile-money dataset of $6{,}362{,}620$
transactions with fields: \texttt{step} (time index, 1 step = 1 hour),
\texttt{type} (PAYMENT, TRANSFER, CASH\_OUT, CASH\_IN, DEBIT),
\texttt{amount}, origin/destination accounts
(\texttt{nameOrig}/\texttt{nameDest}), pre/post balances of both
parties, and the label \texttt{isFraud}.

\subsection{Source 2 --- 9 synthetic features (Faker)}
The project requires additional contextual data generated in Python.
We create 9 fields with the Faker library plus business logic,
\emph{injecting a fraud signal} to reflect realistic correlations (e.g.,
fraudulent accounts tend to be younger and use unfamiliar devices). A
representative snippet is shown in Listing~\ref{lst:gen}.

\begin{lstlisting}[style=pythonstyle, caption={Generating
transaction-level features conditioned on the label (excerpt).},
label={lst:gen}]
fraud = df["isFraud"].to_numpy(bool)
# new device: higher probability if fraud
p_dev = np.where(fraud, 0.55, 0.08)
df["is_new_device"] = (rng.random(n) < p_dev).astype("int8")
# failed payment attempts (Poisson)
lam = np.where(fraud, 2.2, 0.12)
df["failed_payment_attempts"] = rng.poisson(lam).astype("int16")
\end{lstlisting}

All 9 fields are documented in a \emph{data dictionary}
(Table~\ref{tab:datadict}): column name, type, unit, valid range, and
generation logic.

\begin{table}[H]
\centering
\caption{Data dictionary of the 9 synthetic features (source 2).}
\label{tab:datadict}
\small
\begin{tabular}{lll}
\toprule
Field & Type / Unit & Generation logic (summary) \\
\midrule
\texttt{account\_age\_days}          & int32 / days   & Gamma; fraud accounts younger (mean $\sim$44 vs $\sim$500) \\
\texttt{home\_country}               & category       & Uniform draw from a Faker country pool \\
\texttt{device\_fingerprint}         & SHA1           & Sample from a pool of 5{,}000 hashes \\
\texttt{hour\_of\_day}               & int16 / hour   & \texttt{step} \% 24 \\
\texttt{shipping\_billing\_mismatch} & 0/1            & Bernoulli $p{=}0.40$ (fraud) / $0.03$ \\
\texttt{is\_new\_device}             & 0/1            & Bernoulli $p{=}0.55$ (fraud) / $0.08$ \\
\texttt{failed\_payment\_attempts}   & int16 / count  & Poisson $\lambda{=}2.2$ (fraud) / $0.12$ \\
\texttt{ip\_billing\_mismatch}       & 0/1            & Bernoulli $p{=}0.50$ (fraud) / $0.02$ \\
\texttt{ip\_billing\_distance\_km}   & float / km     & Uniform(500,8000) if mismatch, else (0,80) \\
\bottomrule
\end{tabular}
\end{table}

\paragraph{Label-leakage warning.}
Because the fields above are drawn \emph{conditioned on}
\texttt{isFraud}, their predictive power is partly \emph{by design}.
This is why, throughout the project, we compare two feature spaces
(Base vs Full, Section~\ref{sec:feat}) and interpret results cautiously
(Section~\ref{subsec:limit}).

% =========================
\section{Exploratory Data Analysis (Module 2)}
\label{sec:eda}

\subsection{Class imbalance}
The data is severely imbalanced: only $0.129\%$ of transactions are
fraudulent. Hence \textbf{Accuracy is inappropriate}; we prioritize
PR-AUC, Recall, and F1.

\subsection{Fraud by transaction type}
Fraud occurs \textbf{only} in \texttt{TRANSFER} and \texttt{CASH\_OUT}
(Table~\ref{tab:bytype}). This is a key finding: we can pre-filter by
transaction type to raise the positive density.

\begin{table}[H]
\centering
\caption{Fraud distribution by transaction type.}
\label{tab:bytype}
\begin{tabular}{lrrr}
\toprule
Type & \# Fraud & Fraud rate & \# Transactions \\
\midrule
CASH\_OUT & 4{,}116 & 0.184\% & 2{,}237{,}500 \\
TRANSFER  & 4{,}097 & 0.769\% & 532{,}909 \\
CASH\_IN  & 0       & 0\%     & 1{,}399{,}284 \\
PAYMENT   & 0       & 0\%     & 2{,}151{,}495 \\
DEBIT     & 0       & 0\%     & 41{,}432 \\
\bottomrule
\end{tabular}
\end{table}

\subsection{Balance features}
Many fraudulent transactions drain the origin account. We build two
\emph{balance-error} features to capture this anomaly:
\begin{align}
\texttt{errorBalanceOrig} &= \texttt{newbalanceOrig} + \texttt{amount} - \texttt{oldbalanceOrg}, \label{eq:errOrig}\\
\texttt{errorBalanceDest} &= \texttt{oldbalanceDest} + \texttt{amount} - \texttt{newbalanceDest}. \label{eq:errDest}
\end{align}
These separate fraud from legitimate transactions very well and later
dominate the model's feature importance
(Section~\ref{subsec:importance}).

\subsection{Fraud patterns in synthetic signals}
The fraud rate spikes when risk flags equal 1 (new device, address
mismatch, IP mismatch), and \texttt{account\_age\_days} /
\texttt{ip\_billing\_distance\_km} separate clearly by label ---
exactly as designed in Section~\ref{sec:data} (hence the leakage
caveat). In contrast, \texttt{hour\_of\_day} is derived from
\texttt{step} and reflects a \emph{genuine} temporal pattern of PaySim,
with no leakage.

% =========================
\section{Data Cleaning (Module 3)}
\label{sec:clean}

We clean the data under a \emph{before/after documentation} principle,
recording the number of fraud rows removed at every step
(Table~\ref{tab:cleanlog}).

\begin{table}[H]
\centering
\caption{Data-cleaning log.}
\label{tab:cleanlog}
\begin{tabular}{lrr}
\toprule
Step & Rows removed & Fraud removed \\
\midrule
Duplicate records (by original columns)             & 0  & 0 \\
Invalid values (\texttt{amount}$\le$0 / negative balance) & 16 & 16 \\
\texttt{amount} outliers (KEPT, not clipped)        & 0  & 0 \\
\midrule
\textbf{Total} (6{,}362{,}620 $\to$ 6{,}362{,}604)  & \textbf{16} & \textbf{16} \\
\bottomrule
\end{tabular}
\end{table}

\paragraph{Documented decisions.}
\begin{itemize}
    \item \textbf{No missing values or duplicates} --- PaySim is
    simulated, high-quality data.
    \item \textbf{Outliers in \texttt{amount} are kept.} The high-value
    group ($5.31\%$ of the data under the $1.5\times$IQR rule) has a
    \emph{fraud density of $1.14\%$} --- about $8.8\times$ the global
    rate ($0.129\%$). Clipping the tail would delete rare positives;
    here an ``outlier'' is \emph{signal}, not noise.
    \item \textbf{Merchant quirk.} Merchant destination accounts
    (\texttt{nameDest} starting with \texttt{M}, $33.8\%$ of rows)
    always have zero destination balance by PaySim design, and no fraud
    targets a merchant. This is a data characteristic, not an error.
\end{itemize}
Note: all $16$ removed rows were fraud ($16/8{,}213 \approx 0.19\%$) ---
a reminder that ``label-blind'' cleaning can still touch the rare class;
the impact is negligible.

% =========================
\section{Feature Engineering (Module 4)}
\label{sec:feat}

\subsection{Type filtering and two feature spaces}
Based on EDA, we keep only \texttt{TRANSFER} and \texttt{CASH\_OUT},
yielding $2{,}770{,}393$ samples with a positive density raised to
$0.296\%$ (from $0.129\%$). To measure the contribution of the synthetic
data precisely, we build \textbf{two feature spaces}:
\begin{itemize}
    \item \textbf{Base} --- original PaySim transaction features plus
    ``real'' signals derived from the transaction sequence (balance
    errors, velocity). \emph{No label leakage.}
    \item \textbf{Full} --- Base $+$ 7 synthetic risk signals.
\end{itemize}
Every algorithm is trained on \emph{both} to quantify the uplift from
source 2 and to interpret it alongside the leakage caveat.

\subsection{Velocity features, encoding, and scaling}
As required by the project, we add \texttt{txn\_count\_acct} (cumulative
per-account transaction count, causal) and \texttt{time\_since\_last\_txn}
(gap to the previous transaction). The \texttt{type} variable is
binary-encoded as \texttt{type\_TRANSFER}; for Logistic Regression the
numeric features are standardized with \texttt{StandardScaler}.

\paragraph{Practical observation.} On PaySim, accounts almost never
repeat: \texttt{txn\_count\_acct} has mean $\approx 1.0$ and
\texttt{time\_since\_last\_txn} is mostly $-1$. These features are built
\emph{as required} but carry very little signal (see
Section~\ref{subsec:importance}); on a real platform with returning
customers they would matter.

\subsection{Handling class imbalance}
We address imbalance with \textbf{class weights} in all three models
(\texttt{class\_weight="balanced"} for Logistic Regression;
\texttt{scale\_pos\_weight}$\approx$$337$ for XGBoost;
\texttt{class\_weight="balanced\_subsample"} for Random Forest),
combined with the type filtering that already raised positive density.

\subsection{Feature selection}
We rank features by \emph{mutual information} (MI) on a subsample. A
methodological point: \emph{univariate} MI ranks
\texttt{errorBalanceOrig} rather low, \textbf{contrary} to its
dominant importance in tree models --- because its strength lies in
nonlinear interactions. We therefore \textbf{keep all features} for the
tree models (no harm) rather than pruning by univariate MI, only noting
the near-uninformative group (\texttt{newbalanceOrig}).

% =========================
\section{Model Development and Evaluation (Module 5)}
\label{sec:model}

\subsection{Setup}
We split train/test $70/30$, stratified by label (train $1{,}939{,}275$
samples / $5{,}738$ fraud; test $831{,}118$ / $2{,}459$). We train
\textbf{three algorithms} $\times$ \textbf{two feature spaces}: Logistic
Regression, Random Forest, XGBoost. The primary metric is
\textbf{PR-AUC} (suited to imbalanced data), with ROC-AUC, Recall,
Precision, and F1 at the default threshold $0.5$.

\subsection{Comparison of results}
Table~\ref{tab:models} summarizes the 6 models.

\begin{table}[H]
\centering
\caption{Comparison of 6 models on the test set ($831{,}118$ samples,
$2{,}459$ fraud). Threshold $0.5$.}
\label{tab:models}
\small
\begin{tabular}{lccccc}
\toprule
Model & PR-AUC & ROC-AUC & Recall & Precision & F1 \\
\midrule
LogReg $\cdot$ Base & 0.5932 & 0.9782 & 0.9081 & 0.0433 & 0.0827 \\
LogReg $\cdot$ Full & 0.9704 & 0.9999 & 0.9947 & 0.3348 & 0.5010 \\
RandomForest $\cdot$ Base & \textbf{0.9974} & 0.9989 & 0.9967 & 0.9996 & 0.9982 \\
RandomForest $\cdot$ Full & 0.9996 & 1.0000 & 0.9963 & 1.0000 & 0.9982 \\
XGBoost $\cdot$ Base & 0.9944 & 0.9991 & 0.9923 & 0.9204 & 0.9550 \\
XGBoost $\cdot$ Full & \textbf{0.9997} & 1.0000 & 0.9976 & 0.9927 & 0.9951 \\
\bottomrule
\end{tabular}
\end{table}

Key observations:
\begin{itemize}
    \item \textbf{Random Forest $\cdot$ Base is the best honest model}:
    PR-AUC $0.997$ with only \textbf{1 false positive} out of
    $828{,}659$ legitimate transactions (8 fraud missed) --- using
    \emph{no synthetic features at all}.
    \item \textbf{Synthetic features help the linear model greatly}:
    Logistic Regression Base reaches only PR-AUC $0.59$ (precision
    $0.04$, i.e., $48{,}850$ false positives), but Full rises to $0.97$.
    \item On the Base set, XGBoost already reaches $0.994$ --- nearly as
    good as Full --- showing that PaySim is almost separable using
    balance features alone.
\end{itemize}

\subsection{Feature importance}
\label{subsec:importance}
Table~\ref{tab:importance} lists XGBoost-Full importances.
\texttt{errorBalanceOrig} dominates ($0.537$), matching the EDA; notably
the synthetic flags have \emph{low} importance and velocity is $0$ ---
meaning the tree model relies mainly on \emph{real} features, which
tempers (but does not eliminate) the leakage concern.

\begin{table}[H]
\centering
\caption{Feature importance (XGBoost Full, top 10 and tail group).}
\label{tab:importance}
\small
\begin{tabular}{lr@{\hskip 2cm}lr}
\toprule
Feature & Importance & Feature & Importance \\
\midrule
\texttt{errorBalanceOrig}       & 0.537 & \texttt{amount}                 & 0.013 \\
\texttt{account\_age\_days}     & 0.157 & \texttt{shipping\_billing\_mism.} & 0.010 \\
\texttt{newbalanceOrig}         & 0.121 & \texttt{is\_new\_device}        & 0.007 \\
\texttt{failed\_payment\_attempts} & 0.079 & \texttt{oldbalanceDest}      & 0.002 \\
\texttt{ip\_billing\_mismatch}  & 0.028 & \texttt{txn\_count\_acct}       & 0.000 \\
\texttt{type\_TRANSFER}         & 0.016 & \texttt{time\_since\_last\_txn} & 0.000 \\
\bottomrule
\end{tabular}
\end{table}

\subsection{Cost-based metric and operating-threshold selection}
\label{subsec:cost}
The default threshold $0.5$ is not business-optimal. We turn the
Precision/Recall trade-off into \textbf{money}: each missed fraud (FN)
loses exactly the \emph{transaction amount}; each false alarm (FP) costs
a fixed friction $C_{FP}=10$. The total cost at threshold $t$ is
\begin{equation}
\mathrm{Cost}(t) = \!\!\sum_{i:\, y_i=1,\, \hat p_i < t} \!\! \texttt{amount}_i
\;+\; C_{FP}\cdot \bigl|\{i:\, y_i=0,\, \hat p_i \ge t\}\bigr|.
\label{eq:cost}
\end{equation}
Sweeping the threshold on XGBoost Base (the honest, deployment-suited
model) gives an \textbf{optimal threshold $t^{*}=0.21$}
(Table~\ref{tab:cost}), reducing total cost by $\sim$$16.6\%$ versus
$0.5$.

\begin{table}[H]
\centering
\caption{Operating-threshold comparison (XGBoost Base, test set).}
\label{tab:cost}
\begin{tabular}{lrrr}
\toprule
Threshold & Total cost & Recall & FP \\
\midrule
$0.50$ (default) & 1{,}628{,}601 & 0.992 & 211 \\
$0.21$ (optimal) & \textbf{1{,}358{,}351} & \textbf{0.997} & 260 \\
\bottomrule
\end{tabular}
\end{table}

\paragraph{Financial KPI at $t^{*}=0.21$.} The model stops
$\approx 100\%$ of fraud money ($3.6139$B of $3.6152$B on the test set),
letting only $\sim$$1.36$M through, with an \textbf{FPR $=0.031\%$}
($260$ legitimate transactions wrongly blocked) --- a well-balanced
operating point between loss and friction.

\subsection{Limitations and honest interpretation}
\label{subsec:limit}
The near-perfect scores (RF Base PR-AUC $0.997$, 1 FP) should
\textbf{not} be read as excellent capability, but as a sign that the
data is \emph{too easy}:
\begin{itemize}
    \item PaySim simulates fraud with a \emph{near-deterministic balance
    signature} (fraud drains the origin account), so all three tree
    models separate it using real features alone.
    \item The synthetic features are generated from the label, so the
    Full set carries \textbf{leakage}; Full's scores should be seen as
    an \emph{optimistic upper bound}, while the Base set is closer to
    real-world difficulty.
\end{itemize}
We therefore present \emph{both} feature spaces and recommend
\textbf{Random Forest Base} as the best honest model, and
\textbf{XGBoost Base} as the deployment choice (compact, fast, no
leakage).

% =========================
\section{Conclusion}
\label{sec:conclusion}

We built an end-to-end fraud-detection pipeline on PaySim combined with
synthetic data, covering all five modules: business understanding and
KPIs, EDA, documented cleaning, feature engineering with imbalance
handling and feature selection, and training/evaluation of six models
with a cost-based metric. Technically, XGBoost/Random Forest reach
PR-AUC $>0.99$ and the operating threshold $0.21$ stops nearly all fraud
money at an FPR of $0.031\%$. Methodologically, the core contribution is
the \emph{Base-vs-Full comparison design} that exposes the label-leakage
risk, and the conclusion that near-perfect scores reflect the easiness
of the data rather than genuine generalization --- a necessary honest
perspective for reporting.

Future work (Modules 6--8): package the model as an API and review UI,
monitor data drift, and translate findings into risk-policy
recommendations.

\end{document}
